% Solutions to the exercises of Chapter 17.
\begin{solutions}
  \item \ifsolutions
        With $R = \Exp(\delta\theta_g)\hat R \approx (I + \sx{\delta\theta_g})\hat R$ and $a_m - b_a - n_a = \hat a - \delta b_a - n_a$, substituting into
        $\dot v = R(a_m - b_a - n_a) + g_w$, subtracting $\dot{\hat v} = \hat R\hat a + g_w$ and dropping second-order products gives
        \[
          \dot{\delta v} = \sx{\delta\theta_g}\hat R\hat a - \hat R\,\delta b_a - \hat R\,n_a = -\sx{\hat R\hat a}\,\delta\theta_g - \hat R\,\delta b_a - \hat R\,n_a .
        \]
        The local result is $\dot{\delta v} = -\hat R\sx{\hat a}\delta\theta - \hat R\,\delta b_a - \hat R\,n_a$; the two agree through $\delta\theta_g = \hat R\,\delta\theta$ and
        $\sx{\hat R\hat a} = \hat R\sx{\hat a}\hat R\T$. In the velocity row both forms contain $\hat R$, because the velocity error lives in the world frame. The difference
        is in the attitude row: the global error obeys $\dot{\delta\theta}_g = -\hat R\,\delta b_g - \hat R\,n_g$ (\cref{ex:global-attitude}), which depends on the heading through $\hat R$,
        while the local error obeys $\dot{\delta\theta} = -\sx{\hat\omega}\delta\theta - \delta b_g - n_g$, which does not. With the local perturbation the attitude-and-bias
        part of the error model is independent of where the vehicle points (\cref{sec:boxplus}).
        \else\withheldsolution\fi

  \item \ifsolutions
        The attitude row of $F_c$ is $[\,0\ \ 0\ \ {-\sx{\hat\omega}}\ \ 0\ \ {-I}\,]$ in the order $(\delta p, \delta v, \delta\theta, \delta b_a, \delta b_g)$: its $\delta p$ and $\delta v$ blocks are zero, so
        $\delta\theta$ evolves independently of the position and velocity errors. This is the cascade of the kinematics: the attitude is obtained by
        integrating the gyroscope alone ($\dot q = \tfrac12q\qprod[0, \hat\omega]$ involves neither $p$ nor $v$), the velocity by integrating the specific force rotated
        by the attitude, and the position by integrating the velocity. Errors flow downstream only, $\delta b_g \to \delta\theta \to \delta v \to \delta p$ and $\delta b_a \to \delta v$, never
        upstream, so the blocks above the cascade are zero.
        \else\withheldsolution\fi

  \item \ifsolutions
        With the noise ordering $w = (n_a, n_g, n_{b_a}, n_{b_g})$ the noise-input matrix has the blocks $-\hat R$ (velocity, from $n_a$), $-I$ (attitude, from $n_g$), and $I$, $I$
        (the two biases). Since $\hat R\hat R\T = I$,
        \[
          Q_d \approx G_cQ_cG_c\T\Dt = \diag\big(0,\ \sigma_a^2I,\ \sigma_g^2I,\ \sigma_{b_a}^2I,\ \sigma_{b_g}^2I\big)\Dt .
        \]
        Units: $\sigma_a^2\Dt$ is $(\mathrm{m/s^2/\sqrt{Hz}})^2\cdot\mathrm s = \mathrm{m^2/s^2}$, a velocity variance; $\sigma_g^2\Dt$ is $\mathrm{rad^2}$, an attitude variance;
        $\sigma_{b_a}^2\Dt$ is $(\mathrm{m/s^3/\sqrt{Hz}})^2\cdot\mathrm s = \mathrm{m^2/s^4}$, the variance of an acceleration bias; and $\sigma_{b_g}^2\Dt$ is $\mathrm{rad^2/s^2}$. The position block is
        zero to first order: noise reaches the position only through the velocity, contributing $\sigma_a^2\Dt^3/3$ per step at second order
        (\cref{ex:double-integrator-qd}).
        \else\withheldsolution\fi

  \item For a level, stationary IMU each horizontal axis obeys the one-axis chain of \cref{ex:state-models}: $\dot{\delta p} = \delta v$,
        $\dot{\delta v} = g\,\delta\theta - \delta b_a - n_a$, $\dot{\delta\theta} = -\delta b_g - n_g$, with random-walk biases. (Your full $F_c$ reduces to two copies of this chain in
        that configuration, which is a good test of it.) Propagate $P \gets \Phi P\Phi\T + Q_d$ with $\Phi = I + F\Dt$ at 200~Hz from $P = 0$ and plot $\sqrt{P_{pp}}$; it
        follows the analytic total of \eqref{eq:growth-laws}. For the Monte Carlo points, simulate $M = 400$ independent realizations of the same
        discretized equations, drawing $n_a$ and $n_g$ with standard deviations $\sigma/\sqrt\Dt$ and the bias increments with $\sigma_b\sqrt\Dt$
        (\cref{sec:ct-noise}), and record the sample standard deviation of $\delta p$ at logarithmically spaced times. With 400 runs the points scatter by
        about $\pm7\%$ around the curve. The code is \code{fig\_cov\_growth} in \code{scripts/make\_figures.py}.

  \item With no new information the world-frame attitude covariance $\hat RP_{\theta\theta}\hat R\T$ does not change; only the body axes in which $P_{\theta\theta}$ is
        expressed turn (\cref{rem:cov-rotates}). After a quarter turn about $z$, $\hat R(t) = \hat R(0)R_z(90^\circ)$ and
        \[
          P_{\theta\theta}(t) = R_z(90^\circ)\T\,\diag(10^{-4}, 10^{-6}, 10^{-6})\,R_z(90^\circ) = \diag(10^{-6}, 10^{-4}, 10^{-6}):
        \]
        the large and a small diagonal entry have swapped. Along the way, at $45^\circ$, the matrix has equal diagonal entries $5.05\times10^{-5}$ and an
        off-diagonal entry of $\pm4.95\times10^{-5}$. Individual diagonal entries are not monotonic even without noise; the trace and the eigenvalues are constant.
        A test that requires each diagonal entry of $P$ to grow during propagation is therefore wrong (Milestone 5).

  \item A 10\% error in $\sigma_a$ is a 21\% (or, for $-10\%$, 19\%) error in the predicted variance. With $M$ runs, $(M - 1)s^2/\sigma^2 \sim \chi^2_{M-1}$ under a correct prediction,
        and the test rejects when $s^2/\sigma^2$ leaves $[\chi^2_{0.025}(M - 1), \chi^2_{0.975}(M - 1)]/(M - 1)$. The probability of rejecting when the true ratio is $1.21$ (or
        $0.81$) is the power of the test; computing it from the $\chi^2$ distribution gives:
        \begin{center}\small
        \begin{tabular}{@{}l c c c c@{}}
        \toprule
        runs $M$ & 100 & 200 & 500 & 700 \\
        \midrule
        power against $\sigma_a$ 10\% too small (ratio 1.21) & 0.30 & 0.50 & 0.86 & 0.95 \\
        power against $\sigma_a$ 10\% too large (ratio 0.81) & 0.29 & 0.53 & 0.91 & 0.98 \\
        \bottomrule
        \end{tabular}
        \end{center}
        About 600--700 runs are needed to detect a 10\% error with 95\% probability (by the normal approximation, the interval half-width
        $1.96\sqrt{2/M}$ plus $1.645$ times the alternative's spread must be below $0.21$). Two practical remedies: test several independent quantities at
        once (the three axes, several times along the trajectory), which multiplies the effective number of samples, and use larger errors when the
        purpose is only to show that the test can fail.
\end{solutions}
